\documentclass[12pt]{article}
\usepackage[utf8]{inputenc}
\usepackage[T1]{fontenc}
\usepackage{hyperref}
\title{Waja Language Resources for K-12 Curriculum Development}
\author{Compiled by Claude Sonnet 5.5 (claude-sonnet-5-5) for Andr\'as Kornai}
\date{}
\begin{document}
\maketitle

\section{Introduction}
Waja (ISO 639-3 wja; Glottocode waja1259) is a Niger-Congo (Tula--Waja) language of northeastern Nigeria, spoken mainly in Balanga Local Government Area of Gombe State. It is also called Wiyaa or Wuya. Sources give differing speaker counts (about 60,000 in older figures, higher in more recent ones). Very little educational material appears to exist. This report lists only what turned up in web searches.

\section{Primary Educational Resources}
No primers, textbooks or school curricula in Waja were found in search results. The language is not known to be taught in schools. A curriculum would have to be built from scratch, probably starting with the Bible-translation materials below as a source of written text.

\section{Linguistic and Reference Materials}
\begin{itemize}
\item Overview of the language and its classification: \url{https://en.wikipedia.org/wiki/Waja_language}
\item A recent journal article on suffixation in Waja plural formation (Gombe-based authors): \url{https://www.jolan.com.ng/index.php/home/article/view/412}
\item Search results also pointed to an unpublished Waja--English--Hausa dictionary project (University of Maiduguri, 2013) and to Kleinewillinghöfer's 1982 \emph{Untersuchungen zum Waja}. No URL was seen for either, and their availability is unconfirmed.
\end{itemize}

\section{Religious and Community Resources}
\begin{itemize}
\item Waja Bible text on eBible.org, reported as a partial New Testament with a Creative Commons licence: \url{https://ebible.org/Bible/details.php?id=wja}
\item Joshua Project language page, which lists Bible portions and a Jesus Film dubbing: \url{https://joshuaproject.net/index.php/languages/wja}
\end{itemize}
No community organisations were identified in the search results.

\section{Digital and Audio Resources}
No dedicated audio corpus, software or dictionary app was found. The Jesus Film dubbing is mentioned on the Joshua Project page above. The eBible text could serve as a digital corpus.

\section{Curriculum Development Context}
The evidence found is thin: a partial scripture text, a few academic papers and one reference page. There is no established orthography textbook, no school programme and no confirmed dictionary. Any K-12 effort would need community and linguist partnerships, an orthography check against the Bible texts, and new materials. Hausa and English are the likely school languages in the area.

\end{document}